\section*{Capítulo \ref{ch:b1:alcohol-activation} --- Los alcoholes: activar el grupo OH}

\begin{solution}{exo:b1:alcohol-activation:1}
\ce{2CH3CH2OH + 2Na -> 2CH3CH2O- + 2Na+ + H2};
\ce{CH3CH2OH + NaH -> CH3CH2O- + Na+ + H2};
\ce{CH3CH2OH + OH- <=> CH3CH2O- + H2O}, $K = 10^{14.0 - 15.9} = 10^{-1.9}$:
parcial.
\end{solution}

\begin{solution}{exo:b1:alcohol-activation:2}
1-Metoxipropano; metoxietano; etoxibenceno (el fenóxido, una base más
débil que un \omterm{def:b1:alcohol-activation:alkoxide}{alcóxido}, sigue siendo un buen \omterm{def:b1:electronic-effects:nucleophile}{nucleófilo}).
\end{solution}

\begin{solution}{exo:b1:alcohol-activation:3}
La tosilación en piridina da \ce{CH3CH2CH2OTs}; después,
\[ \ce{CH3CH2CH2OTs + I- -> CH3CH2CH2I + TsO-}. \] Directamente, el yoduro
tendría que expulsar \ce{OH-}: no hay reacción. El \omterm{def:b1:alcohol-activation:sulfonate}{tosilato} tiene un \omterm{def:b1:electronic-effects:nucleophile}{grupo
saliente} excelente, y las condiciones no son ni ácidas ni fuertemente
básicas.
\end{solution}

\begin{solution}{exo:b1:alcohol-activation:4}
Butan-2-ol: but-1-eno, \cip{Z}- y \cip{E}-but-2-eno; mayoritario,
\cip{E}-but-2-eno. 2-Metilpentan-2-ol: 2-metilpent-2-eno (mayoritario,
trisustituido) y 2-metilpent-1-eno.
\end{solution}

\begin{solution}{exo:b1:alcohol-activation:5}
\ce{(CH3)2CHCH2-O-CH2CH3}: 2-metilpropan-1-olato de sodio con
bromoetano, o etóxido de sodio con 1-bromo-2-metilpropano. Ambos haluros
son primarios, pero el 1-bromo-2-metilpropano está ramificado junto al
carbono que reacciona, lo que frena la SN2 y deja eliminar a la \omterm{def:b1:predominant-reaction:strong-acid}{base
fuerte} (2-metilpropeno). La primera elección es mejor.
\end{solution}

\begin{solution}{exo:b1:alcohol-activation:6}
El \omterm{def:b1:alcohol-activation:sulfonate}{tosilato} conserva la \omterm{def:b1:stereochemistry-in-depth:isomers}{configuración} \cip{R}; la SN2 con etóxido la
invierte: \cip{S}-2-etoxioctano (el OEt tiene la primera prioridad, como
la tenía el OTs). Con ácido sulfúrico y calor: eliminación a octenos
(sobre todo \cip{E}-oct-2-eno), a través de un \omterm{def:b1:electronic-effects:intermediates}{carbocatión} secundario, con
pérdida de la estereoquímica.
\end{solution}

\begin{solution}{exo:b1:alcohol-activation:7}
Protonación del OH; pérdida de agua hasta el \omterm{def:b1:electronic-effects:intermediates}{carbocatión} terciario
(lenta); captura por \ce{Cl-}: 2-cloro-2-metilpropano. Es rápida porque el
\omterm{def:b1:electronic-effects:intermediates}{carbocatión} terciario es estable y el cloruro, insoluble en el ácido, se
separa.
\end{solution}

\begin{solution}{exo:b1:alcohol-activation:8}
\ce{CH3CH2OH + H+ -> CH3CH2OH2+}; un segundo etanol ataca el carbono por
detrás mientras sale el agua (SN2): \ce{(CH3CH2)2OH+}, que pierde un
protón: etoxietano. En competencia: una base (etanol, \ce{HSO4-}) toma un
protón $\beta$ del \ce{CH3CH2OH2+} mientras sale el agua (E2): eteno.
\end{solution}

\begin{solution}{exo:b1:alcohol-activation:9}
El 2-metilpropan-2-ol (terciario), que da 2-cloro-2-metilpropano,
insoluble.
\end{solution}

\begin{solution}{exo:b1:alcohol-activation:10}
La protonación y la pérdida de agua dan un \omterm{def:b1:electronic-effects:intermediates}{carbocatión} secundario contiguo
a un carbono cuaternario. Un grupo metilo de ese carbono migra con su par
de enlace al carbono catiónico: la carga positiva queda entonces sobre un
\omterm{def:b1:electronic-effects:carbon-class}{carbono terciario}, más estable. La pérdida de un protón de un carbono
vecino da el 2,3-dimetilbut-2-eno tetrasustituido, el alqueno más
estable.
\end{solution}

\begin{solution}{exo:b1:alcohol-activation:11}
A partir del \cip{R}-butan-2-ol: \omterm{def:b1:alcohol-activation:sulfonate}{tosilato} (\cip{R}) y, después, metóxido
de sodio (SN2, inversión): \cip{S}-2-metoxibutano. A partir del
\cip{S}-butan-2-ol: se forma el \omterm{def:b1:alcohol-activation:alkoxide}{alcóxido} (NaH) y se añade yodometano: el
enlace C--O del carbono estereogénico no se toca, y se obtiene de nuevo
\cip{S}-2-metoxibutano.
\end{solution}

\begin{solution}{exo:b1:alcohol-activation:12}
$Q = [\text{éter}][\ce{H2O}]/[\ce{EtOH}]^2$. Eliminar el agua a medida que
se forma mantiene pequeños $[\ce{H2O}]$ y $Q$: $Q < K$ y la reacción sigue
formando éter hasta que se consume el alcohol.
\end{solution}

\begin{solution}{pb:b1:alcohol-activation:1}
\textbf{1.} El terc-butóxido de sodio con un haluro de metilo; el
metóxido de sodio con un haluro de terc-butilo.
\textbf{2.} La segunda: el metóxido, una \omterm{def:b1:predominant-reaction:strong-acid}{base fuerte}, provoca una
eliminación en el haluro terciario,
\ce{(CH3)3CBr + CH3O- -> (CH3)2C=CH2 + CH3OH + Br-}.
\textbf{3.} \ce{(CH3)3CO- + CH3I -> (CH3)3COCH3 + I-}: el oxígeno del
\omterm{def:b1:alcohol-activation:alkoxide}{alcóxido} ataca el carbono del metilo por detrás mientras sale el yoduro
(SN2).
\textbf{4.} Con sodio (o hidruro de sodio): \ce{2(CH3)3COH + 2Na ->
2(CH3)3CO- + 2Na+ + H2}. El hidróxido es una base demasiado débil:
$K = 10^{14.0 -
19.2} = 10^{-5.2}$.
\textbf{5.} El agua protonaría el \omterm{def:b1:alcohol-activation:alkoxide}{alcóxido} y atacaría el haluro.
\textbf{6.} No: no tiene ningún \omterm{def:b1:stereochemistry-in-depth:stereocentre}{centro estereogénico}.
\textbf{7.} \cip{R}-butan-2-ol + TsCl (piridina) $\to$ \omterm{def:b1:alcohol-activation:sulfonate}{tosilato} de
\cip{R}-butan-2-ilo: retención.
\textbf{8.} \cip{S}-2-metoxibutano, por inversión SN2.
\textbf{9.} El \cip{S}-butan-2-ol.
\textbf{10.} \cip{R}-2-metoxibutano: el enlace C--O del carbono
estereogénico se conserva; solo se forma el O--C(metilo).
\textbf{11.} La E2 (el metóxido es una \omterm{def:b1:predominant-reaction:strong-acid}{base fuerte} y el carbono es
secundario), que da butenos; está limitada por la base pequeña, el \omterm{def:b1:electronic-effects:nucleophile}{grupo
saliente} excelente y la temperatura suave.
\textbf{12.} El alcohol protonado se ioniza a un \omterm{def:b1:electronic-effects:intermediates}{carbocatión} secundario
plano, que el metanol captura por sus dos caras: éter racémico, junto con
butenos.
\textbf{13.} Protonación del OH, pérdida de agua (lenta) hasta el
\omterm{def:b1:electronic-effects:intermediates}{carbocatión} terciario y pérdida de un protón $\beta$.
\textbf{14.} 2-Metilbut-2-eno (mayoritario, trisustituido) y
2-metilbut-1-eno.
\textbf{15.} Un \omterm{def:b1:electronic-effects:intermediates}{carbocatión} terciario está impedido y pierde un protón con
mucha más facilidad de lo que otra molécula de alcohol puede atacarlo.
\textbf{16.} Para eliminar un producto: con $Q$ pequeño, la reacción
continúa, y el alqueno no queda en contacto con el ácido.
\textbf{17.} Una primera barrera alta (pérdida de agua, la etapa
determinante) hasta el \omterm{def:b1:electronic-effects:intermediates}{carbocatión} y, después, una barrera pequeña hasta
el alqueno.
\textbf{18.} No: un \omterm{def:b1:electronic-effects:intermediates}{carbocatión} primario es demasiado inestable; los
alcoholes primarios eliminan por E2 sobre el alcohol protonado, a
temperatura más alta.
\textbf{19.} $10.0/74.0 = \qty{0.135}{mol}$.
\textbf{20.} $1.10 \times 0.135 \times 141.9 = \qty{21.1}{g}$.
\textbf{21.} $0.135 \times 88.0 = \qty{11.9}{g}$.
\textbf{22.} $0.75 \times 11.9 =$ \textbf{\qty{8.9}{g} de MTBE}.
\end{solution}
